\clearpage
\fancyhead[L]{\small BODE. ASTRONOMICAL YEARBOOK.}
\hypertarget{native-635-lower}{}
\begin{center}\small Göttingische gelehrte Anzeigen. No. 2. pp. 17..24. 3 January 1824.\\[3pt]\rule{.72\linewidth}{.3pt}\end{center}

\emph{Astronomical Yearbook for the Year \textup{1826} together with a collection of the latest papers, observations and reports pertaining to the astronomical sciences. Calculated and published by} Dr. J. E. \textsc{Bode}, \emph{Royal Astronomer, etc.} Berlin 1823. Available from the author and on commission through Ferd. Dümmler. 256 pages in octavo, with one copperplate.

This volume opens the second set of fifty volumes in a series that has contributed greatly to the spread of astronomical knowledge in Germany and that we hope will long continue \SourceLeaf{636} under the guidance of its worthy editor. Among the celestial events of more general interest in 1826, we single out a total lunar eclipse (on 14 November), a partial solar eclipse (on 29 November) and an occultation of Saturn by the Moon (on 16 February).

The papers and reports in the appendix take up 168 pages. First, some results from the triangulation in the Kingdom of Hanover, by Court Councillor \textsc{Gauss}. This account first gives the geographical positions of 22 selected points as derived in the autumn of 1822 from the preliminary tie to \textsc{von Zach}'s baseline. We may add here that the tie since made to the baseline measured in Holstein by Mr. \textsc{Schumacher} does not alter these results appreciably. It merely reduces the more northerly latitudes by one tenth of an arcsecond. The three towers in Hildesheim that appear unnamed in this list can now be identified. The first is the Andreas tower, the second the Michaelis tower and the third the Jacobi tower. Next comes a summary of the terrestrial refractions inferred from reciprocal zenith distances along 28 lines of the triangulation network. The mean ratio of the curvature of the terrestrial arc to the total observed refraction is 1 to 0.1306. The measurements added in 1823 make the result a little smaller still, namely 1 to 0.1278. This result is considerably smaller than the ratio customarily assumed up to now, when half the refraction (usually, though not quite correctly, called simply refraction) was taken to be 0.08 times the curvature of the terrestrial arc. Yet this determination, in which light from a heliotrope served as the target in almost every observation and which includes so many long distances (the longest is $11\tfrac{1}{2}$ Meilen), seems to deserve more confidence than others in which sighting terrestrial targets in itself did not permit such accuracy in the observed zenith distances and in which, moreover, the distances were usually much shorter. These observations likewise show that the anomalies are always greatest at short distances, especially in the morning, around noon and in the early afternoon. According to the same observations, on sunny summer days in flat country and over shorter distances, light passing close above the ground usually undergoes negative refraction at these times. --- Contributions to the determination of geographical longitudes from observed solar eclipses and stellar occultations, by Prof. \textsc{Wurm}. --- Observations of planetary oppositions, eclipses of Jupiter's satellites and stellar occultations, made at the observatory in Wilna by Mr. \textsc{Sniadecky}. --- Various observations by Mr. \textsc{Rümker} in Paramatta. The most important of them, namely his observations of \textsc{Enke}'s comet in June 1822, have already been reported in no. 26 of this journal for 1823 [p. 442 of this volume]. --- The interesting article that follows, by Dr. \textsc{Olbers}, touches on the most sublime aspect of astronomy, namely the question of whether the universe has bounds or not. Aware of our inability to lift the veil that conceals this mystery, we are nevertheless glad to clear up doubts of the kind that might seem to put the assumption that the universe is infinite in conflict with the rest of our knowledge. One such doubt, at least about the infinite number of suns, is raised by the consideration that, if this number is assumed to be infinite, \emph{every} direction extending from the eye would meet a solar body somewhere, nearer or farther away, and the whole vault of the heavens would therefore appear in an unbroken blaze of sunlight. This inference is in itself perfectly correct. Yet a uniform brilliance would be possible only if space were absolutely transparent, so that light would lose none of its intensity in passing through it. But this assumption is wholly unproved and indeed highly improbable in itself. Dr. \textsc{Olbers} now shows that only a \SourceLeaf{637} comparatively slight degree of opacity need be assumed to remove the apparent contradiction between experience and an infinite number of suns. If the light is weakened by only $\frac{1}{800}$ on its way from Sirius to us, then a hollow sphere with a radius of 30000 times the distance to Sirius, with us at its centre and its entire inner surface radiating the brilliance of the sun, would appear to us only $\frac{1}{700000}$ as bright as the background sky on a clear night with a full moon and would therefore surely make no perceptible impression on our eyes. --- Catalogue of fixed stars by Mr. \textsc{Pond}, communicated by Prof. \textsc{Tralles}. The meaning of this catalogue and of the conclusion drawn from it is not entirely clear to us. --- Further reports on \textsc{Pons}'s (\textsc{Enke}'s) comet by Prof. \textsc{Enke}. The new observations of this celestial body in 1822 confirm the remarkable phenomenon already evident from the earlier observations, that alongside the changes in the orbital period that can be explained by planetary perturbations there is also a progressive decrease in that period, now amounting to a little over three hours from one revolution to the next. In this we probably have the first certain evidence of the actual existence of a resistance that comets experience as they move through space. It is to be hoped that this celestial body can be observed again in Europe as soon as its next approach to the sun in August 1825. At least the work of the calculator, who is as skilful as he is indefatigable, will leave no uncertainty at all as to where to look for it. --- Astronomical observations at the Prag observatory by Messrs. \textsc{David} and \textsc{Bittner}. Similar observations in Kremsmünster by Mr. \textsc{Derflinger} and in Prag by Mr. \textsc{Hallaschka}. --- Original observations of the third comet of 1822 by Dr. \textsc{Olbers}. --- Some mechanical investigations into the origins of comet tails by Dr. \textsc{Lehmann}. It should not surprise us that all attempts so far to explain this puzzling subject satisfactorily have met with no substantial success. We entirely lack observations precise enough to yield exact results on which to base rigorous calculations. We know little for certain about the nature of the forces at work, and treating the effects of forces assumed hypothetically with mathematical rigour presents very great difficulties. Mr. \textsc{Lehmann}'s attempt to explain comet tails solely by known forces, without invoking repulsive forces, still stays far too close to the surface of the problem (as do all the other attempted explanations) and adheres far too closely to inadequate notions in general for it to be possible, on this basis, even to judge whether an explanation along these lines is possible. Nevertheless, it is not without ingenious ideas, and it cannot be denied that the former criticism also applies to the attempts made so far using repulsive forces, whose necessity cannot yet be regarded as settled. --- Astronomical observations at the royal observatory in Berlin in 1822. --- A higher-magnification eyepiece attachment for achromatic telescopes, invented by Mr. W. \textsc{Kitchiner} in London. As far as one can gather from the very incomplete and unclear description, the invention consists merely in mounting the eyepiece lenses so that they can move relative to each other, thereby obtaining variable magnification. --- Observations and elements of the comet of September 1822 by Mr. \textsc{Rümker} in Paramatta. These observations extend to around the middle of November, considerably later than the European observations, and therefore deserve to be combined with them. Also, an observation by the same observer of the transit of Mercury on 5 November 1822. --- Determinations of geographical position in the Altmark and other borderlands by Mr. \textsc{Stöpel}, music director in Tangermünde. This valuable catalogue contains 134 determinations of position based on trigonometric measurements linked to the triangle side from Magdeburg to Hagelsberg in the great triangulation of \textsc{von Müffling}. We should have liked Mr. \textsc{Stöpel} also to have told us something about the means he used and his \SourceLeaf{638} method of observation, instead of merely referring us on that point to a provincial paper unknown elsewhere. --- Observation of the occultation of a star by Mars, by Prof. \textsc{Tralles}, all the more remarkable given the rarity of such phenomena. After emersion, the star did not become visible again until it was more than one planetary diameter away. It was very faint at first but then brightened rapidly. The telescope had an aperture of $3\tfrac{1}{2}$ inches. Mr. \textsc{Tralles} infers from this the existence of a substantial Martian atmosphere. --- On the photosphere of the planets discovered by Mr. G. R. \textsc{Pastorf}, by Mr. \textsc{Ritz}. Mr. \textsc{Ritz} explains this phenomenon by a double reflection between the objective lenses. The reviewer has always held the same view and has satisfied himself, by calculation using the actual dimensions of \textsc{Frauenhofer}'s telescopes, that the reflection in question must cause those telescopes to produce a phenomenon like the one observed. An essential condition for its occurrence is that the curvatures of the two surfaces facing one another must be \emph{nearly} equal, as is indeed the case in the design of \textsc{Frauenhofer}'s objectives. --- Remarks on the preceding subject by Judicial Commission Councillor \textsc{Kunovsky}. Mr. \textsc{Kunovsky} approves of the explanation offered by Mr. \textsc{Ritz}, whose paper appears to have been shared with him for inspection before printing. The reviewer admits, however, that Mr. \textsc{Kunovsky} seems to him to have misunderstood the spirit of the explanation, since almost everything he cites as confirmation, if one were entitled to regard it as established fact, would contradict that explanation or at least require a quite different explanation of its own. Mr. \textsc{Kunovsky} claims to have noticed the phenomenon with \emph{all} the telescopes he tested, although the English objectives are known to be constructed quite differently from \textsc{Frauenhofer}'s. He finds that with different telescopes the sphere of light is nearly proportional to the aperture of the objective. In the spirit of the explanation above, this holds only for \emph{one and the same} telescope insofar as the objective is allowed a larger aperture in one instance and a smaller one in another. If different telescopes had unequal dimensions that were proportional to one another in \emph{all} respects, the sphere of light would have to appear exactly the same size. Finally, when using much higher magnifications, Mr. \textsc{Kunovsky} saw Jupiter's sphere of light become much fainter, as it must, whereas the sphere of light of $\alpha$ in the Lyre is said to have lost scarcely any of its brightness at higher magnification, which again would be incompatible with the explanation above. The test of the purity of the glass in the objectives suggested by Mr. \textsc{Kunovsky} is, incidentally, very useful. --- A letter from Prof. \textsc{Bürg} supplies some results of his more recent investigations into lunar theory. --- Fragments toward the interpretation of Aratus, by Prof. \textsc{Schaubach}. --- Observation of the total lunar eclipse on 26 January 1823 and of the occultation of a star, by Mr. \textsc{Rümker} in Paramatta. --- Some observations at the Dorpat Observatory. --- Similar observations at the observatory in Krakau, by Prof. \textsc{Leski}. --- On astronomical refraction, by Prof. \textsc{Bessel}. --- Astronomical news from the Rev. \textsc{Luthmer} in Hannover. \textsc{Herschel}'s epitaph. --- New elements of Juno's orbit, by Prof. \textsc{Nicolai}. These are the very valuable results of a laborious calculation, based on quadratures, of the perturbations Jupiter has exerted on Juno's motion since 1804, and of a discussion of 15 oppositions observed since the discovery. The calculation represents those oppositions with great accuracy. Mr. \textsc{Nicolai} is inclined to attribute the small remaining discrepancies to a suspected inadequacy of \textsc{Newton}'s law of gravitation. Specifically, he considers it possible that at the same distance Jupiter exerts a different attraction on the Sun than on Juno. Such a hypothesis, which others too have already put forward, cannot of course be declared absolutely impossible. Yet it is highly improbable if only because it has been established as a fact that the Sun's attraction on all primary and secondary planets \SourceLeaf{639} is equally great, just as the most diverse bodies on the Earth's surface gravitate toward the Earth equally. In the reviewer's opinion, therefore, one should resort to an assumption so contrary to analogy only once decisive evidence exists for it, which is by no means the case so far. --- Stellar occultations observed at the Vienna Observatory by Prof. \textsc{Littrow}. --- List of 795 double stars from Mr. \textsc{Struve}'s Astr. Observations, volume 3. --- There are also a number of interesting things among the shorter astronomical notices and remarks, although there is still room for objections to several articles. For example, the reviewer does not quite see how the fact that the observed orbital periods of \textsc{Halley}'s comet differ much more \emph{from one another} than those of \textsc{Enke}'s comet is supposed to demonstrate the great advances in the art of observation and in calculation, undeniable though those advances in themselves are. Even the fact that the difference between the \emph{calculated} orbital period and the observed one was much greater for \textsc{Halley}'s comet than for \textsc{Enke}'s does little in itself to demonstrate advances in calculation, since the circumstances were so very different for the two comets.
